\section{Conclusion}
\label{sec:conclusion}

We characterised a structural bias in the dominant evaluation protocol for cross-lingual multimodal generation. On the full $n=2142$ jointly-judged MTVQA validation set across three open VLMs, the Latin/non-Latin artifact share spans 121 pp (Qwen $+66\%$, InternVL $+8\%$, Phi-3.5-vision $-57\%$), and strict-EM and semantic-EM reverse the Qwen-vs-InternVL ranking in $100\%$ of 2000 paired bootstrap resamples; output-style verbosity (non-Latin/Latin character ratio) rank-orders the three architectures identically to artifact share. Across 13 plausible interventions on Qwen-7B, aggregate $\Delta$strict-EM moves up to $+4.2$ pp while aggregate $\Delta$semantic-EM-C is statistically null; the largest-lifting SFT recipe yields $\Delta$strict $+3.64$ pp [$+0.06$, $+7.23$] with $\Delta$sem-C $-0.28$ pp [$-4.98$, $+4.42$], and per-language point estimates show direction-of-effect anti-correlation on non-Latin scripts. The choice of multilingual VLM from a strict-EM leaderboard is not the choice of stronger model; it is the choice of architecture whose output-style prior matches gold-string convention, with a sign and magnitude that depend on the architecture. We propose a nine-item reporting protocol (R1--R9) deployable at single-digit dollar cost. Future work: replicate to DocVQA multilingual and free-form benchmarks; extend to closed-source baselines (Claude, GPT-4o vision, Gemini-Pro); design SFT and preference-learning recipes whose loss is computed against semantic equivalence rather than gold-string exact match.
